\section{A stable family and its thermodynamic functions}
\label{f:section}

We first construct a bounded repulsive core which permits clusters of at
most four particles at very short separations.  Its quadratic control of
local particle numbers will be uniform over all attraction ranges used
later.  Throughout the proof, all position integrals use Lebesgue measure
in one fixed microscopic length unit, and all finite volumes have free
boundaries.

\subsection{Packing and the microscopic core}

\begin{lemma}[Packing density]\label{f:packing}
Let $M_L$ be the maximum cardinality of a subset of $[0,L]^3$ whose
distinct points have mutual distances at least one.  Then
\[
  p=\lim_{L\to\infty}\frac{M_L}{L^3}
\]
exists and belongs to $(0,\infty)$.
\end{lemma}

\begin{proof}
Disjoint balls of radius $1/3$ about the points lie in
$[-1/3,L+1/3]^3$, so $M_L\le C(L+1)^3$.  This also bounds the
possible cardinalities; compactness of the configuration space for each
cardinality shows that the maximum is attained.  The integer lattice
gives a positive lower asymptotic density.

Fix $s>0$ and repeat a maximizing configuration in $[0,s]^3$ by
translations in $(s+1)\mathbb Z^3$.  Points in distinct translates are
separated by at least one: in a coordinate in which the translates
differ, their coordinate separation is at least $(s+1)-s=1$.
Taking all complete translates contained in $[0,L]^3$ gives
\[
 \liminf_{L\to\infty}\frac{M_L}{L^3}
 \ge \frac{M_s}{(s+1)^3}.
\]
Letting $s\to\infty$ along a sequence realizing the upper asymptotic
density makes the right side tend to that upper density.  The lower
and upper limits therefore agree.
\end{proof}

Set $A=3/(2p)$ and $\Delta=10^{-3}$.  The parameter $t>1$ below is a
construction parameter, distinct from inverse temperature $\beta$.
We shall eventually choose one sufficiently large $t$ and keep it
fixed.  Put
\[
 l_t=\ee^{-t},\qquad I_t=[l_t,(1+\Delta)l_t],
 \qquad 2(1+\Delta)l_t<1.
\]
Define $\psi_t:[0,\infty)\to[0,t^2]$ by
\begin{equation}\label{f:core}
 \begin{aligned}
 \psi_t(r)&=t^2
   &&\text{on }[0,l_t]\cup[(1+\Delta)l_t,1],\\
 \psi_t(r)&=0
   &&\text{on }[(1+\Delta/3)l_t,(1+2\Delta/3)l_t]
                    \cup[1+1/t,\infty),
 \end{aligned}
\end{equation}
and make it affine on each of the three intervening intervals.  Thus
$\psi_t$ is continuous, including at zero, nonnegative, and supported
in $[0,2]$.

Fix a radial nonnegative $C^1$ function $w$ on $\R^3$, supported in the
closed unit ball, with $\int_{\R^3}w(x)\,\dd x=1$.  For example, one
may normalize $((1-\abs{x}^2)_+)^2$.  Write
$w_R(x)=R^{-3}w(x/R)$.  When its argument is a displacement in $\R^3$,
we also write $\psi_t$ for the radial lift $x\mapsto\psi_t(\abs{x})$.
The family of potentials used below is
\begin{equation}\label{f:family}
 \varphi(x)=\psi_t(\abs{x})-S(x),\qquad
 S(x)=t\sum_j\alpha_j w_{R_j}(x),\qquad
 R_j\ge1,\quad \alpha_j\ge0,\quad \sum_j\alpha_j\le A+1.
\end{equation}
The index set may be empty, finite, or countable.  Since
$\norm{w_{R_j}}_\infty\le\norm{w}_\infty$, the sum converges
uniformly in the countable case.  Every such potential is therefore
bounded, continuous, radial, and nonpositive when $\abs{x}\ge2$.

The following lemma combines an elementary geometric obstruction with
an extremal graph estimate.  The weighted edge maximization in its
proof is the Motzkin--Straus argument for Tur\'an's bound
\citep[Theorem~1]{MotzkinStraus1965}.  We include it to make the
quadratic core estimate explicit.

\begin{lemma}[The shell graph]\label{f:shell-graph}
For any finite configuration in $\R^3$, join two particles when their
distance belongs to $I_t$.  This graph has no clique of size five.
On any $m$ vertices it has at most $3m^2/8$ edges.
\end{lemma}

\begin{proof}
If five points formed a clique, rescale their distances by $l_t^{-1}$
and take the four difference vectors from one of them.  Their Gram
matrix differs entrywise by at most $2\Delta+\Delta^2$ from the matrix
with diagonal one and off-diagonal entries $1/2$.  For a diagonal entry
this follows from $1\le r^2\le(1+\Delta)^2$; for an off-diagonal entry
use polarization,
$\ip{v_i}{v_j}=(\abs{v_i}^2+\abs{v_j}^2-\abs{v_i-v_j}^2)/2$.
The reference matrix has smallest eigenvalue $1/2$.  The perturbation
has operator norm at most $4(2\Delta+\Delta^2)<1/2$, so the Gram matrix
would be positive definite and have rank four.  Four vectors in
$\R^3$ cannot have this Gram matrix.

For the edge bound, maximize $\sum_{\{i,j\}\in E}u_i u_j$ over
$u_i\ge0$, $\sum_i u_i=1$, choosing a maximizer with the smallest
possible support.  If two nonadjacent vertices have positive weights,
the objective is affine in their two weights with their sum fixed.
Putting their combined weight on one of the vertices cannot decrease
the objective, contradicting minimality of the support.  The support
is therefore a clique of some size $r\le4$.  Its objective is
\[
 \frac12\left(1-\sum_i u_i^2\right)
 \le\frac12(1-1/r)\le\frac38.
\]
For $m>0$, apply this bound to equal weights $1/m$.  The empty graph
case is immediate.
\end{proof}

\subsection{Quadratic control of cell counts}

We now convert the exclusion of five-particle shell cliques into a
lower bound that controls arbitrary particle numbers, including
coincident particles.  This is a concrete superstable estimate in the
sense of \citet{Ruelle1970}; the constants will be uniform over all
attraction ranges.

Partition space into half-open cubes $C_i$ of side $b=1/4$, with any
fixed translation of the grid.  When working in a bounded region,
replace each cell by its intersection with that region and omit empty
cells.  Let $n_i$ be the number of particles in cell $i$.  If
\[
 J_{ij}=\sup_{x\in C_i,\,y\in C_j}S(x-y),
\]
then $J$ is symmetric and
\begin{equation}\label{f:row-bound}
 \sup_i\sum_jJ_{ij}\le M:=c_b t(A+1),
 \qquad c_b=\norm{w}_\infty(2/b+5)^3.
\end{equation}
Indeed, a cell interacts through $w_R$ with at most $(2R/b+5)^3$
cells, each contributing at most $R^{-3}\norm{w}_\infty$.  Since
$R\ge1$, their row sum is at most $c_b$.  Summing the coefficients
proves \eqref{f:row-bound}, also for a countable sum and cropped cells.

We shall need the same estimates after removing some interactions.
More precisely, in any bounded region allow symmetric measurable
pair interactions
\begin{equation}\label{f:modified-energy}
 v(x,y)=c(x,y)\psi_t(\abs{x-y})-K(x,y),
 \qquad 0\le c\le1,\quad K\ge0,
\end{equation}
where $c(x,y)=1$ for points in the same small cell, $K$ is bounded,
and the symmetric matrix
$J_{ij}=\sup_{x\in C_i,y\in C_j}K(x,y)$ has row sums at most $M$.
The coefficients are fixed functions of the positions, independent of
the particle configuration.  This includes partial deletion of any
attractive pairs, partial deletion of core pairs between distinct
cells, and their linear interpolations whenever they share the stated
row bound.  No translation invariance is required for these finite
volume estimates.  Write $U^v(X)=\sum_{r<s}v(x_r,x_s)$, and write
$U^\varphi$ for a member of \eqref{f:family}.

\begin{proposition}[Common stability]\label{f:stability}
Set
\begin{equation}\label{f:stability-constants}
 D_0=t^2/8,\qquad B_0=t^2/2.
\end{equation}
For every configuration, including coincident particles,
\begin{equation}\label{f:cell-stability}
 U^{\psi_t}\ge D_0\sum_i n_i^2-B_0N,
 \qquad
 U^v\ge(D_0-M/2)\sum_i n_i^2-B_0N.
\end{equation}
Choose $t$ sufficiently large that
\begin{equation}\label{f:t-requirements}
 D_0\ge8M,\qquad ct>4(A+1),\qquad c=1/3456.
\end{equation}
Then, in any axis-parallel cube of side $L\ge1$, with $V=L^3$,
\begin{equation}\label{f:density-stability}
 U^v\ge ct^2\frac{N^2}{V}-\frac{t^2}{2}N.
\end{equation}
In particular all members of \eqref{f:family} are stable with the
same constant $B_0$.
\end{proposition}

\begin{proof}
The diameter $\sqrt3 b$ of a small cell is less than one.  By
\Cref{f:shell-graph}, among $m$ particles in one cell at least
$m(m-1)/2-3m^2/8=m^2/8-m/2$ pairs have distance outside $I_t$.
Each of these pairs contributes $t^2$, while all remaining core
contributions are nonnegative.  Summing over cells proves the first
bound, and also proves it when core terms between cells are deleted.
For the attractive contribution, including extra diagonal self terms
only enlarges the bound, and symmetry gives
\[
 \sum_{r<s}K(x_r,x_s)
 \le\frac12\sum_{i,j}J_{ij}n_i n_j
 \le\frac14\sum_{i,j}J_{ij}(n_i^2+n_j^2)
 \le\frac M2\sum_i n_i^2.
\]
This proves \eqref{f:cell-stability}.

There are at most $(L/b+2)^3\le216V$ nonempty cells.  Hence
$\sum_i n_i^2\ge N^2/(216V)$.  Under \eqref{f:t-requirements},
$D_0-M/2\ge t^2/16$, giving \eqref{f:density-stability} with the stated
$c$.  The requirements on $t$ are compatible: $M$ grows linearly in
$t$, whereas $D_0$ grows quadratically.  The constants $c$ and $c_b$
are independent of $t$ and of the attraction ranges.
\end{proof}

The two margins in \eqref{f:t-requirements} will be used separately.
The inequality $D_0\ge8M$ makes the influence of neighboring cell counts
small compared with the repulsion within one cell; this will yield
volume-independent occupation moments.  The inequality $ct>4(A+1)$
keeps the density bound coercive even after adding the quadratic
attractive weight in the next grand-sum estimate.

\subsection{Canonical limits and pressure duality}

The local quadratic bound controls configurations of high density.
We next use it to establish the thermodynamic functions required by
the attraction transform.  The proof applies directly to the infinite
attraction sum: separated boxes interact only attractively, so no
finite-range approximation is needed to construct these limits.

For $\Lambda_L=[0,L]^3$, define the canonical partition functions
and their normalization by
\begin{equation}\label{f:canonical-definition}
 \begin{aligned}
 Z^\varphi_{L,n}(\beta)
 &=\frac1{n!}\int_{\Lambda_L^n}\exp(-\beta U^\varphi(X))\,\dd X,
 & Z^\varphi_{L,0}&=1,\\
 H^\varphi_{L,n}(\beta)&=-\frac1{tL^3}\log Z^\varphi_{L,n}(\beta).
 \end{aligned}
\end{equation}

\begin{lemma}[Finite volume canonical bounds]\label{f:canonical-bounds}
There is an absolute constant $C<\infty$, independent of $t$ and all
attraction ranges, such that for every admissible $t$, $L\ge1$,
$n\ge0$, $\beta>0$, and $d_L=n/L^3$,
\begin{equation}\label{f:finite-bounds}
 \beta t(cd_L^2-d_L/2)+\frac{d_L\log d_L-d_L}{t}
 \le H^\varphi_{L,n}(\beta)
 \le \beta Ct d_L^2+\frac{d_L\log d_L}{t},
\end{equation}
where $0\log0=0$.
\end{lemma}

\begin{proof}
The lower bound follows by bounding the integrand from above using
\eqref{f:density-stability} and using $n!\ge(n/\ee)^n$ for $n\ge1$.
For the other bound, let $X_1,\ldots,X_n$ be independent uniform
points in $\Lambda_L$.  Since $U^\varphi\le U^{\psi_t}$ and
$\int_{\R^3}\psi_t(\abs{x})\,\dd x\le t^2\abs{B(0,2)}$,
\[
 \mathbb E U^\varphi
 \le\frac{n(n-1)}{2L^3}\int_{\R^3}\psi_t(\abs{x})\,\dd x
 \le Ct^2n^2/L^3.
\]
Jensen's inequality gives
$\log Z^\varphi_{L,n}\ge n\log L^3-\log n!-\beta Ct^2n^2/L^3$.
Now use $n!\le n^n$.  At $n=0$ both bounds are equalities.
\end{proof}

The following bound controls the entire large-particle-number tail.
The additional quadratic factor is included because it will also
occur in the range comparison.

\begin{lemma}[Grand sum tails]\label{f:grand-tail}
Fix an admissible $t$ and a compact $\mathcal D\subset\Omega$, where
$\Omega=(0,\infty)\times\R$ has coordinates $(\beta,h)$.
There are constants $D\ge1$ and $a_0>0$, depending only on $t$ and
$\mathcal D$, such that, for every $\varphi$ in \eqref{f:family},
$L\ge1$, and $0\le\zeta\le2(A+1)$,
\begin{equation}\label{f:grand-tail-bound}
 \sum_{n>DL^3}\ee^{thn}Z^\varphi_{L,n}(\beta)
          \exp\left(\frac{\beta t\zeta n^2}{2L^3}\right)
 \le\frac{\exp[-a_0(\lfloor DL^3\rfloor+1)]}{1-\ee^{-a_0}}
 \qquad ((\beta,h)\in\mathcal D).
\end{equation}
In particular the tail tends uniformly to zero, and is at most one
for all sufficiently large $L$.  These constants require no local
occupation estimate.
\end{lemma}

\begin{proof}
Enclose $\mathcal D$ in $\beta\in[\beta_-,\beta_+]$ and
$\abs{h}\le H$.  Put $V=L^3$.  By \eqref{f:finite-bounds}, the logarithm
of the summand at $n\ge1$ is at most
\[
 -\beta t(ct-\zeta/2)\frac{n^2}{V}
 +(th+\beta t^2/2+1)n-n\log(n/V).
\]
The choice \eqref{f:t-requirements} implies
$ct-\zeta/2\ge ct/2$.  Set
$\kappa=\beta_-ct^2/2$ and $B=tH+\beta_+t^2/2+1$.
For $n/V\ge1$ the displayed expression is at most
$-\kappa n^2/V+Bn$.  Choose $D\ge\max(1,2B/\kappa)$.
For $n>DV$ this is at most
$-\kappa n^2/(2V)\le-(\kappa D/2)n$.
Summing this geometric majorant proves the assertion with
$a_0=\kappa D/2$.
\end{proof}

\begin{proposition}[Thermodynamic limits and duality]
\label{f:thermodynamics}
For every member $\varphi$ of \eqref{f:family}, including countably
infinite sums, and every $\beta>0$, $d\ge0$, the limit
\begin{equation}\label{f:canonical-limit}
 H_\varphi(\beta,d)
 =\lim_{\substack{L\to\infty\\ n/L^3\to d}}
      H^\varphi_{L,n}(\beta)
\end{equation}
exists along every such sequence of real sides and integer counts.
It is finite, convex and continuous in $d\in[0,\infty)$, and concave
in $\beta>0$ for each fixed $d$.  The bounds \eqref{f:finite-bounds}
hold with $H_\varphi(\beta,d)$ and $d$ in place of their finite volume
counterparts.

The grand partition function and pressure satisfy
\begin{equation}\label{f:pressure-definition}
 \Xi^\varphi_L(\beta,h)
 =\sum_{n\ge0}\ee^{thn}Z^\varphi_{L,n}(\beta),
 \qquad
 Q_\varphi(\beta,h)
 =\lim_{L\to\infty}\frac{\log\Xi^\varphi_L(\beta,h)}{tL^3}.
\end{equation}
The pressure is finite, nonnegative, and jointly convex on $\Omega$.
For $\beta>0$,
\begin{equation}\label{f:duality}
 \begin{split}
 Q_\varphi(\beta,h)
   &=\sup_{d\ge0}\{hd-H_\varphi(\beta,d)\},\\
 H_\varphi(\beta,d)
   &=\sup_{h\in\R}\{hd-Q_\varphi(\beta,h)\},\qquad d>0.
 \end{split}
\end{equation}
For fixed admissible $t$, pressures of the whole family are uniformly
bounded on each compact subset of $\Omega$.
\end{proposition}

\begin{proof}
We first prove \eqref{f:canonical-limit}, without using a pressure
limit or an approximation by finite-range potentials.  At $d=0$,
both bounds in \eqref{f:finite-bounds} tend to zero.  Fix $d>0$ and
$\beta>0$.  Let $H_*$ be the lower limit of
$H^\varphi_{L,n}(\beta)$ as $L\to\infty$ and $n/L^3\to d$, and choose
$s_j\to\infty$ and integers $n_j$ with
\[
 n_j/s_j^3\longrightarrow d,\qquad
 H^\varphi_{s_j,n_j}(\beta)\longrightarrow H_*.
\]
The finite volume bounds make $H_*$ finite.  Consider any other
sequence $L\to\infty$, $n/L^3\to d$.  For each fixed, sufficiently
large $j$, place
$k=\lfloor L/(s_j+2)\rfloor^3$ cubes of side $s_j$ in $\Lambda_L$,
separated by gaps of width at least two.  For any nonnegative integers
$m_1,\ldots,m_k$ summing to $n$, restriction to configurations with
these box counts gives
\begin{equation}\label{f:canonical-product}
 Z^\varphi_{L,n}(\beta)
 \ge\prod_{i=1}^k Z^\varphi_{s_j,m_i}(\beta).
\end{equation}
Indeed, interactions between different boxes are nonpositive.  There
are $n!/\prod_i m_i!$ disjoint assignments of particle labels to the
boxes, so the global factor $1/n!$ leaves exactly the product of
their factors $1/m_i!$.

We retain the favorable count $n_j$ in most boxes and correct the
total particle number in the remaining boxes.  Only a vanishing fraction
of boxes will change, and \eqref{f:finite-bounds} will bound their
canonical costs uniformly.  To make this exact for every target count, put
\[
 q_j=\lfloor(4d+4)s_j^3\rfloor,
 \qquad c_d=\min(d/2,1).
\]
For all sufficiently large $j$,
$n_j\ge c_ds_j^3$ and $q_j-n_j\ge c_ds_j^3$.
Start with $n_j$ particles in every box.  If
$E=n-kn_j\ge0$, fill boxes successively to $q_j$, with at most one
last box given an intermediate integer count.  If $E<0$, empty boxes
successively, again with at most one last intermediate count.  This
changes at most $r$ boxes, where
\begin{equation}\label{f:integer-adjustment}
 \frac rk\le\frac1{c_d}
       \abs{\frac{n}{ks_j^3}-\frac{n_j}{s_j^3}}+\frac1k.
\end{equation}
For fixed large $j$ and then sufficiently large $L$, the target lies
between $0$ and $kq_j$: its mean count divided by $s_j^3$ tends to
$d(1+2/s_j)^3<q_j/s_j^3$.  Thus the adjustment never exhausts the
available boxes.  The quotient and remainder in division of $E$ by
$q_j-n_j$, or of $-E$ by $n_j$, give the asserted integer counts.

All counts used have density in $[0,4d+4]$.  By
\eqref{f:finite-bounds}, their normalized canonical costs have absolute
value at most some $K<\infty$, uniformly for large $j$.  Taking
logarithms in \eqref{f:canonical-product} consequently yields
\[
 H^\varphi_{L,n}(\beta)
 \le\frac{ks_j^3}{L^3}
       \left(H^\varphi_{s_j,n_j}(\beta)+2K\frac rk\right).
\]
First let $L\to\infty$ with $j$ fixed.  The prefactor tends to
$(s_j/(s_j+2))^3$, and the upper limit of $r/k$ is at most
\[
 c_d^{-1}\abs{d(1+2/s_j)^3-n_j/s_j^3}.
\]
Then let $j\to\infty$.  This error tends to zero and the preceding
upper bound tends to $H_*$.  Every sequence has lower limit at least
$H_*$ by its definition; thus every sequence has the same limit.
This proves \eqref{f:canonical-limit} with the required order of limits.

For density convexity, fix $d_1,d_2\ge0$ and $\theta\in[0,1]$.
Choose $s_j\to\infty$, integers $m_{r,j}$ with
$m_{r,j}/s_j^3\to d_r$ for $r=1,2$, and integers $q_j\to\infty$.
In a cube of side $q_j(s_j+2)$, use $q_j^3$ separated boxes of side
$s_j$, assigning count $m_{1,j}$ to $\lfloor\theta q_j^3\rfloor$
of them and count $m_{2,j}$ to the rest.  The total density tends to
$\theta d_1+(1-\theta)d_2$.  The product inequality and the already
proved all-sequence limit, applied directly to this diagonal sequence,
give
\[
 H_\varphi(\beta,\theta d_1+(1-\theta)d_2)
 \le\theta H_\varphi(\beta,d_1)
       +(1-\theta)H_\varphi(\beta,d_2).
\]
This proof does not presuppose continuity in density.  A finite convex
function is continuous at positive densities: on any interior compact
interval its secant slopes are bounded by secants to two exterior
points.  Continuity at zero follows from \eqref{f:finite-bounds}.
H\"older's inequality makes $\log Z^\varphi_{L,n}(\beta)$ convex in
$\beta$.  Passing to the canonical limit proves the stated concavity.

The grand sums are finite by \Cref{f:grand-tail}, with $\zeta=0$.
Fix $(\beta,h)$.  Choose $D$ there and write $V=L^3$.  Above $DV$ the
entire sum tends to zero.  Among the remaining
$\lfloor DV\rfloor+1$ terms choose a maximizing count $n_L$; the
maximum is at least one because it includes the vacuum.  Therefore
\[
 \frac{\log\Xi^\varphi_L}{tV}
 \le h\frac{n_L}{V}-H^\varphi_{L,n_L}(\beta)
       +\frac{\log(\lfloor DV\rfloor+2)}{tV}
\]
for all sufficiently large $L$.  From any subsequence choose a further
subsequence with $n_L/V\to d\in[0,D]$.  The canonical limit bounds
its upper limit by $hd-H_\varphi(\beta,d)$.  Conversely, for every
fixed $d\ge0$ the single term $n=\lfloor dV\rfloor$ bounds the lower
limit by $hd-H_\varphi(\beta,d)$.  This proves existence of the
pressure and the first duality in \eqref{f:duality}.

The vacuum gives $Q_\varphi\ge0$.  The canonical lower bound and
$d\log d-d\ge-1$ bound the first supremum in \eqref{f:duality} by a
finite quadratic maximum, uniformly over the family on any fixed
parameter compact.  H\"older's inequality on the disjoint union of
all configuration spaces makes each $\log\Xi^\varphi_L$ jointly
convex in $(\beta,h)$; its pointwise limit has the same property.

Finally fix $d>0$.  Convexity of the finite function
$H_\varphi(\beta,\cdot)$ gives a supporting slope $h$ there.  Explicitly,
one may choose $h$ between
\[
 \sup_{0\le v<d}\frac{H_\varphi(\beta,d)-H_\varphi(\beta,v)}{d-v}
 \quad\text{and}\quad
 \inf_{v>d}\frac{H_\varphi(\beta,v)-H_\varphi(\beta,d)}{v-d}.
\]
The left endpoint is at most the right endpoint by convexity, and
both are finite by comparison with secants through fixed points on
either side of $d$.  The resulting supporting inequality holds for
all $v\ge0$, so the first duality gives
$Q_\varphi(\beta,h)=hd-H_\varphi(\beta,d)$.  For all other $h$ the
same first duality gives the opposite variational inequality.  This
proves the second equality in \eqref{f:duality} without a differentiability
or uniqueness assumption.
\end{proof}
